\documentclass[12pt]{article}

\usepackage{amsmath}

\begin{document}

Estoy tratando de resolver el siguiente problema:

Dos cargas puntuales q1 y q2 están separadas una distancia d.

a) Hallar el trabajo que es necesario realizar para traer en forma
cuasiestacionaria otra carga q desde un punto muy alejado hasta el punto
central del segmento que separa a q1 y q2.

b) Analice el resultado si las cargas son de igual valor absoluto y de signo
diferente.  Discuta la relación de los resultados con la dirección del campo
eléctrico (Ayuda: considere la mediatriz del segmento que une ambas cargas y la
irrotacionalidad del campo electrostático).

Mi respueta al punto a) es:

\begin{align*}
    W=\frac{1}{4\pi\epsilon_0}\frac{1}{d/2}(q_1+q_2)
\end{align*}

\noindent y al punto b) es $W=0$ pues si $q_1=-q_2$ y muevo $q_0$ por la
mediatriz la fuerza eléctrica resultante sorbre $q_0$ es siempre ortogonal al
desplazamiento de $q_0$, por lo que el trabajo de esta fuerza es zero.

¿Te parecen correctas mis respuestas? No entiendo porque la ayuda hace
referencia a la irrotacionalida del campo electrostático.

\end{document}
